\subsection{Integración de las Pruebas en el \textit{Pipeline} de Integración Continua}\label{subsec:pruebas-ci}

El valor de una \textit{suite} de pruebas se maximiza cuando su ejecución es automática y obligatoria, de modo que ninguna integración degrade las garantías ya verificadas. El \textit{pipeline} de integración continua en GitHub Actions —descrito en el capítulo de desarrollo (\autoref{subsec:sdlc}) y evidenciado en el de resultados (\autoref{subsec:resultado-d})— se ejecuta de forma obligatoria frente a cada \textit{pull request} dirigida a las ramas protegidas, aplicando dos compuertas de calidad: la validación estática diferencial (\textit{linting} y formato sobre los archivos modificados) y la ejecución automática de la \textit{suite} de reglas de seguridad, para la cual aprovisiona la \textit{Firebase Emulator Suite} y corre el \textit{script} \texttt{test:rules} contra el emulador de Firestore.

Así, la garantía de seguridad deja de depender de una acción manual y se convierte en una compuerta automática del proceso de integración: una \textit{pull request} cuyas modificaciones rompan una regla de seguridad es bloqueada antes de promoverse, evitando que una regresión en el blindaje de los datos alcance \textit{staging} o producción. Mejoras incrementales como el cacheo del emulador y la publicación de artefactos de cobertura se retoman en el apartado de trabajo futuro.
